\begin{thebibliography}{200}  
% 调整所有条目间距
\setlength{\itemsep}{0pt}      % 条目之间的垂直间距
\setlength{\parskip}{0pt}       % 段落间距

\bibitem{ref1}新战略移动机器人产业研究所,CMR产业联盟.具身智能复合移动机器人产业发展蓝皮书（2025版）[R].2025.

\bibitem{ref2}付明玉,白丹,张坦,等.具有未知扰动上界的气垫船全鲁棒滑模控制[J].控制理论与应用,2022,39(06):1130-1138.

\bibitem{ref3}凌明,张志鹏,杨勇.RISC-V嵌入式系统设计[M].北京:机械工业出版社,2025.

\bibitem{ref4}王元慧,张峻恺,吴鹏.基于强化学习自抗扰的气垫船进坞控制策略[J].哈尔滨工程大学学报,2025,43(7):1340-1348.

\bibitem{ref5}肖剑波,陆爱杰,陈婧,等.基于模糊PID控制器的气垫船航向控制系统[J].舰船科学技术,2019,41(9):76-80.

\bibitem{ref6}崔延青.全垫升气垫船航向—横倾解耦控制研究[D].哈尔滨:哈尔滨工程大学,2017.

\bibitem{ref7}刘文浩,仇海鸿,米佳琪.无刷电机与视觉感知结合的智能控制气垫船[EB/OL].国家级大学生创新创业训练计划平台,(2025-07-19)[2025-12-27].

\bibitem{ref8}大连理工大学创新团队.多功能两栖气垫船[EB/OL].挑战杯官方网站,(2011)[2025-12-27].

\bibitem{ref9}Alves J C, Ramos T M, Cruz N A. A reconfigurable computing system for an autonomous sailboat[C]//Oceans 2008. IEEE,2008:1-6.

\bibitem{ref10}Liceaga-Castro E, Liceaga-Castro J. Towards a multivariable automatic control system design for an air cushion vehicle[J].IFAC Proceedings Volumes,1995,28(17):335-340.

\bibitem{ref11}Parameter Estimation and Verification of Unmanned Air Cushion Vehicle (UACV) System[C]//MATEC Web of Conferences.EDP Sciences,2017,97:01069.

\bibitem{ref12}Flex-RV Team.Bendable non-silicon RISC-V microprocessor[J].Nature,2024,634(8033):341-346.

    
\end{thebibliography}